\section{STRIDE Analysis}\label{sec:stride-analysis}

This section aims to further explain the threat mapping presented in Table~\ref{tab:stride-vuln-matrix}. 

\subsection{Spoofing}

\textbf{Password-only:} Password-only systems are the most vulnerable to phishing atttack, since they entirely rely on the user's better judgement, there is no origin binding, or relying on a safe third party. 

\textbf{2FA:} 2FA creates a much stronger barrier against phishing attacks, since it provides another wall of defence, in the form of a second factor. Depending on the type of this factor attackers can use different ways to bypass it mostly via adversary-in-the-middle proxies, or sim swapping.

\textbf{FIDO2/WebAuthn:} FIDO2's origin binding ensures that even if an attacker manages to phish the user, they can't redirect the user to a fake site, since the system would not accept the authentication request from a foreign RP ID.

\subsection{Tampering}

\textbf{Password-only:} Usually password only systems offer no protection against message tampering beyond TLS. If TLS is compromised or misconfigured, an attacker can modify login requests or responses.

\textbf{2FA:} 2FA TOTP systems also rely primarily on TLS for tampering protection. OTP codes can be be intercepted during transmission, for example in an AiTM relay attack (Real Time Phishing), where the attacker uses a proxy to phish the user's password and intercept the OTP code, then uses it to authenticate.

\textbf{FIDO2/WebAuthn:} Since FIDO2 is a digitally signed challenge-response system, modification to the challenge or response invalidates the signature, making it virtually tamper-proof even if TLS is compromised.

\subsection{Repudiation}

\textbf{Password-only:} Repudiation protection in a barebones password system is non-existent, anyone with the password can authenticate, and there is no way to link a particular login to a particular device or authenticator. Even logging can be manipulated since we can not be sure about the user's identity. 

\textbf{2FA:} 2FA can provide some repudiation protection through audit logs of OTP usage, but these logs can modified by attackers who gain high enough privileges. The OTP code itself doesn't provide cryptographic proof of which device generated it, which makes it  even more difficult to prove who actually performed the authentication.

\textbf{FIDO2/WebAuthn:} FIDO2 provides strong non-repudiation through cryptographic signatures. Each authentication is signed by a specific authenticator (identified by credential ID and public key), creating an audit trail that cannot be forged. The signature counter also helps detect credential cloning attempts.

\subsection{Information Disclosure}

\textbf{Password-only:} Password systems are highly vulnerable to information disclosure. Passwords are stored as hashes, but breaches can expose these hashes for offline cracking. More critically, passwords are transmitted during every login, creating multiple opportunities for interception. Password reuse across sites amplifies the impact of any single breach.

\textbf{2FA:} TOTP reduces information disclosure risk compared to passwords. The TOTP secret is stored server-side and never transmitted, but if compromised, it can generate valid codes. OTP codes themselves are short-lived but can be intercepted during transmission. Backup codes, if stored improperly, can also be a disclosure risk.

\textbf{FIDO2/WebAuthn:} FIDO2 minimizes information disclosure risk. No shared secrets are stored or transmitted—only public keys are stored server-side, and private keys never leave the authenticator. Even if the server database is breached, attackers cannot use the public keys to authenticate. The challenge-response mechanism ensures no sensitive data is transmitted during authentication.

\subsection{Denial of Service}

\textbf{Password-only:} Password systems are vulnerable to DoS through brute-force attacks and account lockout mechanisms. Attackers can trigger lockouts by repeatedly attempting incorrect passwords, preventing legitimate users from accessing their accounts. Rate limiting helps but can be bypassed through distributed attacks.

\textbf{2FA:} TOTP systems face DoS risks from both password-based attacks and OTP flooding. Attackers can exhaust backup codes, flood users with OTP requests, or trigger account lockouts. Push-based MFA is particularly vulnerable to push fatigue attacks, where repeated prompts cause users to accidentally approve malicious requests.

\textbf{FIDO2/WebAuthn:} FIDO2 is more resistant to DoS attacks (Single user DoS). Brute-force attacks are infeasible due to public-key cryptography—each attempt requires a valid signature, which cannot be guessed. Account lockouts are less effective as attack vectors since there are no shared secrets to guess. However, users can still be locked out if they lose their authenticator and don't have backup credentials configured.

\subsection{Elevation of Privilege}

\textbf{Password-only:} Password systems are vulnerable to privilege elevation through weak or reused credentials. Administrative accounts with weak passwords can be compromised, and password reuse across systems allows lateral movement. The simplicity of password-based authentication makes it easier for attackers to gain initial access.

\textbf{2FA:} TOTP provides some protection against privilege elevation by requiring an additional factor, but high-privilege accounts remain targets for phishing and OTP relay attacks. If an administrator falls victim to phishing, their elevated privileges can be exploited immediately.

\textbf{FIDO2/WebAuthn:} FIDO2 provides strong protection against privilege elevation. High-privilege accounts tied to FIDO2 credentials require physical access to the authenticator device, making remote compromise significantly more difficult. The origin binding prevents credential theft through phishing, and the cryptographic nature of authentication makes privilege escalation through credential compromise nearly impossible.

\subsection{Summary}

The STRIDE analysis reveals a clear security hierarchy: password-only systems are vulnerable across all threat categories, 2FA provides moderate improvements but remains susceptible to sophisticated attacks, and FIDO2/WebAuthn offers strong protection against most threat vectors. The primary security advantages of FIDO2 stem from its use of public-key cryptography, origin binding, and the elimination of shared secrets, making it resistant to the most common authentication attack vectors identified in Section~\ref{sec:common-vulnerabilities}.


